\section{Il modello a due temperature}
\label{sec:modello}

\daScrivere

% Traccia:
%  - decomposizione dell'energia nei modi (eq. 1-7 del paper) e le due T
%  - variabili conservate: rho_s, E, E_ve (eq. 23); perche' energie e non T
%  - equazione di stato: p = somma delle pressioni parziali (eq. 24)
%  - scambio V-T: Landau-Teller (eq. 8), tau di Millikan-White + Park (eq. 9-17)
%  - scambio V-V: Knab (eq. 18), solo nel caso N2-O2
%  - chimica: legge di azione di massa (eq. 27), Arrhenius (eq. 28),
%    temperatura di Park T^a Tv^(1-a) (eq. 29)
%  - accoppiamento chimica-vibrazione: modello non preferenziale (eq. 30-31)
%  - riepilogo del sistema risolto nel bagno termico (eq. 22, 25, 26)
